\chapter{实验复现与参数配置说明}

为保证本文实验过程具有可复现性，系统在原型验证阶段对设备连接、检测模型、控制策略和目标记忆模块进行了统一配置。实验前需要保证布控球设备与验证主机之间具有稳定网络连接，并完成设备 SDK、播放库接口和检测模型的可用性检查。在此基础上，系统通过实时预览链路获取视频帧，通过检测线程输出人体目标信息，并由控制线程根据目标偏差生成 PTZ 控制指令。

实验阶段采用 YOLOv8n 作为人体检测模型，输入尺寸设置为 384，置信度阈值设置为 0.4。该配置优先考虑实时性，以满足检测结果能够及时反馈至 PTZ 控制模块的需求。人体 ReID 记忆模块采用低维颜色和几何特征进行短时身份关联，匹配阈值设置为 0.86，目标最大保留帧数设置为 120，用于支持目标短时遮挡或离开画面后的重新关联。

PTZ 控制模块采用阈值死区、比例脉冲、指令冷却和指数平滑相结合的控制策略。水平和垂直方向分别设置中心死区，目标中心点在进入死区范围后不再触发云台动作；当偏差超过阈值时，系统根据偏差大小调整单次控制脉冲时长，并通过冷却时间限制连续指令频率。该参数组合能够在原型阶段兼顾响应速度和控制稳定性。

在实验复现过程中，应重点记录视频链路延迟、检测帧率、目标 ID 连续性、PTZ 响应方向和异常退出后的资源释放情况。上述指标能够反映系统从视频输入、目标感知、控制决策到设备执行的完整闭环效果，也为后续嵌入式平台迁移和复杂场景优化提供对比依据。
